% GENERATED FILE. Edit canonical lessons, then run npm run generate:books.
% canonical-source-hash: fnv1a64:7d173f2493b2b80b
% canonical-lessons: ES-C56-dad

\chapter{-dad --- Twelve Hundred More}
\label{ch:friend-dad}
\hlchaptermodality{24}

\begin{chapteropening}
\textbf{By the end of this chapter:} \emph{I can read a -dad word I have never met, and say why it looks less like its English partner.}

Complete ES-C56-dad: convert university, read realidad unseen, and explain the uneven erosion.
\end{chapteropening}

\section[-dad]{-dad — twelve hundred more}

\label{lesson:ES-C56-dad}



\begin{quote}
\emph{-ción} is \emph{-tion}. Here is a second ending, and this one is harder to spot because the two languages wore it down by different amounts.
\end{quote}

\begin{sounds}
\begin{itemize}
\raggedright
  \item \texttt{d-soft} — the \emph{d} between vowels is soft, close to English \emph{th} in \emph{this}: \emph{ciudad} ends nearer \emph{-dath} than \emph{-dad}.
\end{itemize}
\end{sounds}

\begin{grammarlens}[title={Grammar Lens: -dad is -ty}]
\begin{quote}
\textbf{Spanish \emph{-dad} (or \emph{-tad}) is English \emph{-ty}.}
\end{quote}

\begin{quote}
\emph{ciudad} $\to$ city · \emph{libertad} $\to$ liberty · \emph{universidad} $\to$ university \emph{realidad} $\to$ reality · \emph{dificultad} $\to$ difficulty · \emph{sociedad} $\to$ society
\end{quote}

This one does not announce itself the way \emph{-ción} does. \emph{Ciudad} and \emph{city} do not look alike at a glance — but line the endings up and the rule is exact, and once you have seen it you cannot unsee it.

\textbf{Feminine again}, every time: \emph{la ciudad}, \emph{la libertad}. That is now two endings that carry their own gender, which between them cover a large slice of the abstract nouns in the language.
\end{grammarlens}

\begin{cousinweb}
Latin \textbf{-tātem} made a noun meaning \emph{the quality of being} something.

Spanish inherited it and wore it to \textbf{-dad}: the \emph{t} softened to \emph{d} between vowels, exactly as it does in the middle of \emph{ciudad} today. English took it through French, where it had already worn much further — \emph{-tātem} to \emph{-té} to \textbf{-ty}.

So English is holding a shorter stub of the same ending. That is why this pair looks less alike than \emph{-ción}/\emph{-tion}: not a different relationship, just more erosion on the English side.

Which gives you a way to read the ones you have not met. See a Spanish word ending in \emph{-dad} and you are looking at an abstract quality, and English almost certainly has the same word with a smaller tail.
\end{cousinweb}

\subsection*{Your turn}
Say these aloud:

\begin{itemize}
\raggedright
  \item \textquotedblleft{}ciudad, libertad, universidad\textquotedblright{} — the same tail three times
  \item \textquotedblleft{}city, liberty, university\textquotedblright{} straight after, and hear the shorter tail
\end{itemize}

\begin{tcolorbox}[breakable,colback=teal!4,colframe=teal!35!black,title={Before you move on}]
English \emph{-ty} becomes? (\emph{\textbf{-dad}}.) Gender? (\textbf{Feminine}.) And what does a \emph{-dad} word always name? (\textbf{A quality} — the state of being something.) Next: an ending you can put on almost any adjective.
\end{tcolorbox}
